\documentclass{article}

\usepackage{lipsum}
\usepackage{amsmath}
\usepackage{graphicx}
\usepackage{subcaption}
\usepackage{hyperref}
\usepackage{xcolor} % For coloring text
\usepackage{amsmath} 
\usepackage{comment}

\renewcommand{\figurename}{Figure}


% Punkte-Variable definieren
\newcounter{points}
\newcommand{\addpoints}[1]{\addtocounter{points}{#1}}

% Punkte-Umgebung für Aufgaben
\newenvironment{task}[1]
  {
    \par \noindent \textbf{Aufgabe #1} \par
    \addpoints{#1}
  }
  {\par}

\begin{document}

\section*{Obstacle Avoidance}
Fix the range, to 180 Degree, get the middle of each segment with the angle, and put the provided force into it, so smooth out noise.

\begin{figure}[H]
    \centering
    \includegraphics[width=0.5\linewidth]{FixStuff/pictures/ObstacleAvoidance.pdf}
    \caption{}
    \label{fig:bluearrow}
\end{figure}

\section*{Mapping}

Fix noises, sometimes if there's a wall it detects something further between the wall.
-> Done

Fix free space in front of robot, if scan is nan or inf


\end{document}
